\section{例题与分析}
\begin{question}
\begin{example}[指数衰减与指数增长]
某类粒子的衰变常量为 $\lambda>0$，即每个粒子在很短时间 $\dd t$ 内的衰变概率为 $\lambda\dd t$。初始期望数为 $N_0$，求未衰变期望数和半衰期。另讨论理想种群按连续比例率 $\lambda$ 增长的数量。
\end{example}
\fig{c4e1q}{短时间内的减少量与当时数量成正比}{\begin{tikzpicture}
\node[draw](a)at(0,0){$N(t)$};\node[draw](b)at(4.5,0){$N(t)+\dd N$};
\draw[flow](a)--(b)node[midway,above]{$\dd N=-\lambda N\dd t$};\node[below]at(2.25,-.5){$\lambda$ 的单位为时间的倒数};
\end{tikzpicture}}
\end{question}
\begin{solution}\textbf{把相对变化率积分并代入初值}\par\nobreak
图 \ref{c4e1q} 中的局部规律为 $\dot N=-\lambda N$。$N_0>0$ 时分离变量，
\[
 \int_{N_0}^{N(t)}\frac{\dd n}{n}
 =-\lambda\int_0^t\dd s
 \quad\Rightarrow\quad
 \boxed{N=N_0\e^{-\lambda t}}.
\]
左侧积分为 $\ln N-\ln N_0=\ln(N/N_0)$，右侧为 $-\lambda t$。两边取指数，得到 $N/N_0=\e^{-\lambda t}$。

\par\medskip\Needspace{4\baselineskip}\noindent\textbf{由衰减解反求半衰期}\par\nobreak
令 $N=N_0/2$，则 $-\ln2=-\lambda t_{1/2}$，所以 $t_{1/2}=\ln2/\lambda$。
\fig{c4e1s}{衰减与增长由同一初值向不同方向演化}{\begin{tikzpicture}[x=1.3cm,y=1cm]
\draw[axis](0,0)--(3.5,0)node[right]{$\lambda t$};\draw[axis](0,0)--(0,3)node[above]{$N/N_0$};
\draw[blue,domain=0:3,samples=50]plot(\x,{exp(-\x)});\draw[red,domain=0:1,samples=40]plot(\x,{exp(\x)});\node[left]at(0,1){$1$};
\node[right,blue]at(2,.3){衰减};\node[right,red]at(1,2.7){增长};
\end{tikzpicture}}


\par\medskip\Needspace{4\baselineskip}\noindent\textbf{改用增长率并检查零初值}\par\nobreak
理想种群的方程为 $\dot N=+\lambda N$，图 \ref{c4e1s} 的增长支对应 $N=N_0\e^{\lambda t}$。$N_0=0$ 时两方程都有恒零解。连续比例率与“每过一个单位时间增加固定百分数”的离散模型需区分：若离散增长因子为 $1+p$、单位时长为 $\Delta t$，对应连续率为 $\ln(1+p)/\Delta t$。
\end{solution}

\begin{question}
\begin{example}[线性阻力下的滑行]
质量 $m$ 的物体在水平光滑面上仅受阻力 $F=-kv$，$k>0$。初始 $x=0,v=v_0>0$。求位置、速度以及最终趋近的位置，判断是否在有限时间停下。
\end{example}
\fig{c4e2q}{阻力始终与当前速度相反}{\begin{tikzpicture}
\draw[axis](-.5,0)--(6,0)node[right]{$x$};\draw[fill=gray!15](1.8,0)rectangle(2.7,.6);
\draw[force,blue](2.7,.5)--(4,.5)node[above]{$v$};\draw[force,red](1.8,.4)--(.6,.4)node[above]{$kv$};
\node[below]at(0,0){$0$};
\end{tikzpicture}}
\end{question}
\begin{solution}\textbf{在正速度支路上求衰减规律}\par\nobreak
图 \ref{c4e2q} 取向右为正，$m\dot v=-kv$。初速度为正，在速度不为零的运动段可分离为 $\dd v/v=-(k/m)\dd t$。两侧按初值积分，
\[
\int_{v_0}^{v(t)}\frac{\dd u}{u}
=-\frac{k}{m}\int_0^t\dd s,\qquad
[\ln u]_{v_0}^{v(t)}=-\frac{k}{m}[s]_0^t.
\]
代入端点，得到$\ln[v(t)/v_0]=-kt/m$。
定义 $\tau=m/k$，取指数后得速度；再从 $x(0)=0$ 对速度积分，得到
\[
 v=v_0\e^{-t/\tau},\qquad
 x=\int_0^t v_0\e^{-s/\tau}\dd s
 =v_0\tau(1-\e^{-t/\tau}).
\]
\fig{c4e2s}{速度渐近为零，位置渐近到有限值}{\begin{tikzpicture}[x=1.1cm,y=1.8cm]
\draw[axis](0,0)--(4.8,0)node[right]{$t/\tau$};\draw[axis](0,0)--(0,1.4);
\draw[blue,domain=0:4.3,samples=60]plot(\x,{exp(-\x)});\draw[red,domain=0:4.3,samples=60]plot(\x,{1-exp(-\x)});
\draw[guide](0,1)--(4.4,1);\node[right,blue]at(4.4,.15){$v/v_0$};\node[right,red]at(4.4,.85){$x/(v_0\tau)$};
\end{tikzpicture}}



\par\medskip\Needspace{4\baselineskip}\noindent\textbf{求极限距离并判断是否有限时停止}\par\nobreak
位置积分的原函数是 $-v_0\tau\e^{-t/\tau}$，上下限代入后为 $-v_0\tau\e^{-t/\tau}+v_0\tau$，所以得到上式。图 \ref{c4e2s} 显示 $t\to\infty$ 时 $x\to v_0\tau=mv_0/k$；有限 $t$ 时指数始终为正，故没有有限停止时刻。

另在 $v>0$ 的段内，$m v\,\dd v/\dd x=-kv$ 可除以 $v$，变为 $\dd v=-(k/m)\dd x$。取$x(0)=0,v(0)=v_0$，将变量与上下限对应后积分：
\[
\int_{v_0}^{v}\dd u=-\frac{k}{m}\int_0^x\dd \xi,
\qquad
[u]_{v_0}^{v}=-\frac{k}{m}[\xi]_0^x.
\]
所以$v-v_0=-kx/m$，即$v=v_0-kx/m$。它给出同一极限位置，却不能单独判定到达时间；时间结论须由指数解确定。
\end{solution}

\begin{question}
\begin{example}[速度与剩余距离成反比]
追逐者与静止目标初距 $L$，初速 $v_0$，随后速度满足 $v=v_0L/r$，$r$ 为剩余距离。求相遇时间。若两者都按同样规则对称迎面运动，初速各为 $v_0$，相遇时间怎样改变？
\end{example}
\fig{c4e3q}{剩余距离缩短时速度增加}{\begin{tikzpicture}
\draw[axis](0,0)--(6,0);\fill(1,0)circle(2pt)node[below]{追逐者};\fill(5,0)circle(2pt)node[below]{目标};
\draw[flow](1,.3)--(2.2,.3)node[above]{$v_0L/r$};\draw[dim](1,-.8)--(5,-.8)node[midway,below]{$r$};
\end{tikzpicture}}
\end{question}
\begin{solution}\textbf{用剩余距离的减小确定微分方程符号}\par\nobreak
图 \ref{c4e3q} 的距离随时间减小，因此 $\dot r=-v_0L/r$，不是正号。先分离变量$r\,\dd r=-v_0L\,\dd t$。从初距$r(0)=L$开始，
\[
\int_L^{r(t)}\rho\,\dd\rho=-v_0L\int_0^t\dd s,
\qquad
\left[\frac{\rho^2}{2}\right]_L^{r(t)}=-v_0L[s]_0^t.
\]
代入上下限，得到$(r^2-L^2)/2=-v_0Lt$，移项后$r^2=L^2-2v_0Lt$。相遇时 $r=0$，所以 $2v_0Lt_c=L^2$，即 $t_c=L/(2v_0)$。
\fig{c4e3s}{对称迎面运动时，闭合速率为两者速度之和}{\begin{tikzpicture}
\fill(0,0)circle(2pt);\fill(5,0)circle(2pt);\draw[flow](0,.2)--(1.3,.2);\draw[flow](5,.2)--(3.7,.2);
\draw[dim](0,-.5)--(5,-.5)node[midway,below]{$r$};\node[above]at(2.5,.6){$-\dot r=2v_0L/r$};
\end{tikzpicture}}


\par\medskip\Needspace{4\baselineskip}\noindent\textbf{双向接近时重新计算闭合速率}\par\nobreak
图 \ref{c4e3s} 中，两者各自朝对方前进，闭合速率为两个速率之和，所以 $\dot r=-2v_0L/r$。乘以$r\,\dd t$，再使用同一初值$r(0)=L$：
\[
\int_L^{r(t)}\rho\,\dd\rho=-2v_0L\int_0^t\dd s,
\qquad
\left[\frac{\rho^2}{2}\right]_L^{r(t)}=-2v_0L[s]_0^t.
\]
所以$(r^2-L^2)/2=-2v_0Lt$，即$r^2=L^2-4v_0Lt$。令 $r=0$，得到 $t_c=L/(4v_0)$。两结果都是理想速度规律一直外推到 $r=0$ 的极限时间；速度本身在那里发散。若存在速度上限，应在达到上限后换用新规律。
\end{solution}

\begin{question}
\begin{example}[容器泄流]
竖直圆柱容器截面积 $S$，初始液面距底部孔口的高度为 $H$，小孔面积 $S_0\ll S$。液体不可压、上下均受大气压，忽略损耗。求液面距孔口的高度 $h(t)$ 及从 $H$ 降至 $H/2$ 的时间。
\end{example}
\fig{c4e4q}{液面下降与孔口流出由体积守恒联系}{\begin{tikzpicture}
\draw[world](-1,2.7)--(-1,0)--(1,0);\draw[world](1,.25)--(1,2.7);\fill[blue!10](-.98,.02)rectangle(.98,2.1);\draw[blue](-1,2.1)--(1,2.1);
\draw[flow](1,.12)--(2.2,.12)node[right]{$v$};\draw[guide](-1.6,.12)--(-1,.12);\draw[dim](-1.5,.12)--(-1.5,2.1)node[midway,left]{$h$};
\node[right]at(1.1,.5){$S_0$};\node[above]at(0,2.1){截面 $S$};
\end{tikzpicture}}
\end{question}
\begin{solution}\textbf{由孔口出流与容器体积变化列连续方程}\par\nobreak
图 \ref{c4e4q} 中，液面与孔口均受大气压，相同压强项在能量关系中抵消。液面处单位质量的重力势能为 $gh$；孔口高度取零，单位质量动能为 $v^2/2$。因 $S_0\ll S$，液面速度相对孔口可忽略，故 $gh=v^2/2$、$v=\sqrt{2gh}$。

短时间内流出体积为 $S_0v\dd t$，容器内体积从 $Sh$ 变为 $S(h+\dd h)$，所以变化量为 $S\dd h<0$。二者相加应为零，得到
\[
 S\dot h=-S_0\sqrt{2gh},\qquad
 \frac{\dd h}{\sqrt h}=-\frac{S_0}{S}\sqrt{2g}\dd t.
\]
\fig{c4e4s}{根号高度随时间线性下降}{\begin{tikzpicture}
\draw[axis](0,0)--(5,0)node[right]{$t$};\draw[axis](0,0)--(0,2.8)node[above]{$\sqrt h$};
\draw[blue,thick](0,2.3)--(4.3,0);\node[left]at(0,2.3){$\sqrt H$};\node[below]at(4.3,0){$t_{\rm empty}$};
\end{tikzpicture}}



\par\medskip\Needspace{4\baselineskip}\noindent\textbf{带初始液面积分，并在排空处终止支路}\par\nobreak
按 $h:H\to h(t)$、$t:0\to t$ 积分，左侧是 $2(\sqrt h-\sqrt H)$，右侧是 $-(S_0/S)\sqrt{2g}\,t$。除以二并移项，图 \ref{c4e4s} 对应
\[
 \sqrt h=\sqrt H-\frac{S_0}{S}\sqrt{\frac g2}\,t,
\quad
 t_{1/2}=\boxed{\frac S{S_0}\sqrt{\frac{2H}{g}}
 \left(1-\frac1{\sqrt2}\right)}.
\]
为求半高时间，将 $\sqrt h$ 换成 $\sqrt{H/2}=\sqrt H/\sqrt2$，有
$(S_0/S)\sqrt{g/2}\,t_{1/2}=\sqrt H(1-1/\sqrt2)$，除去时间系数便得到上式。排空时同理取 $\sqrt h=0$，有 $t_{\rm empty}=(S/S_0)\sqrt{2H/g}$。高度公式只在 $0\le t\le t_{\rm empty}$ 使用。若把平方后的式子继续延伸，会得到液深重新增加的非物理解。由流量守恒，液面速率为 $(S_0/S)v$，所以被忽略的动能与孔口动能之比为 $(S_0/S)^2$。
\end{solution}

\begin{question}
\begin{example}[等温大气压强]
大气为理想气体，摩尔质量 $\mu$，温度恒为 $T$，重力加速度 $g$ 视为常数。地面压强 $p_0$，求高度 $z$ 处的压强和密度。
\end{example}
\fig{c4e5q}{水平薄层上下表面的压力不同}{\begin{tikzpicture}
\draw[axis](-1,0)--(-1,3)node[above]{$z$};
\draw[fill=blue!10](0,1.1)rectangle(2.7,1.5);
\draw[force](1.35,.2)--(1.35,1.1)node[midway,left]{$p(z)S$};\draw[force](1.35,2.6)--(1.35,1.5)node[midway,right]{$p(z+\dd z)S$};
\node[right]at(2.8,1.3){$\dd z$};\draw[force,red](2.2,1.1)--(2.2,.3)node[midway,right]{$\rho gS\,\dd z$};
\end{tikzpicture}}
\end{question}
\begin{solution}\textbf{对气体薄层列受力平衡并消去密度}\par\nobreak
图 \ref{c4e5q} 薄层质量为 $\rho S\dd z$。竖直平衡给出
$p(z)S-p(z+\dd z)S-\rho gS\dd z=0$，即 $\dd p/\dd z=-\rho g$。
图中的红色向下箭头是薄层自重，不能在平衡式中遗漏。理想气体满足 $pV=\nu RT$，薄层质量为 $\mu\nu$，所以 $\rho=\mu\nu/V=\mu p/(RT)$。代入静水平衡式，得到 $\dd p/\dd z=-\mu gp/(RT)$，再分离变量并用 $p(0)=p_0$ 积分，
\[
 \frac{\dd p}{p}=-\frac{\mu g}{RT}\dd z,\qquad
 \boxed{p=p_0\e^{-z/H_s}},\quad H_s=\frac{RT}{\mu g}.
\]
\fig{c4e5s}{压强和密度具有同一标高}{\begin{tikzpicture}[x=1cm,y=1cm]
\draw[axis](0,0)--(4,0)node[right]{$z/H_s$};\draw[axis](0,0)--(0,2.4)node[above]{$p/p_0$};
\draw[blue,domain=0:3.5,samples=50]plot(\x,{2*exp(-\x)});\node[left]at(0,2){$1$};
\draw[guide](1,0)--(1,.7358)--(0,.7358);\node[below]at(1,0){$1$};\node[left]at(0,.7358){$\e^{-1}$};
\end{tikzpicture}}



\par\medskip\Needspace{4\baselineskip}\noindent\textbf{由地面边界得到压强与密度分布}\par\nobreak
具体积分为
$\int_{p_0}^{p(z)}\dd p/p=-(\mu g/RT)\int_0^z\dd z$，
即 $\ln[p(z)/p_0]=-\mu gz/(RT)$。取指数后得到上式；再代回 $\rho=\mu p/(RT)$，便有
$\rho(z)=\mu p_0\e^{-z/H_s}/(RT)$。
图 \ref{c4e5s} 因而同时表示压强和密度的归一化变化。$H_s$ 的单位为长度，温度越高衰减越慢；使用中要求温度与 $g$ 在该高度范围内近似不变。
\end{solution}

\begin{question}
\begin{example}[恒定切向力驱动的细杆]
均匀细杆质量 $m$、长 $L$，一端绕固定光滑轴 $O$ 自由转动。另一端始终受大小 $F$、垂直于杆的力，力矩方向不变。忽略重力，从静止开始求转角、角速度和端点加速度分量。
\end{example}
\fig{c4e6q}{力随杆转动，保持垂直于杆}{\begin{tikzpicture}
\fill(0,0)circle(2.5pt)node[below]{$O$};\draw[very thick](0,0)--(3,1.5)node[midway,below]{$L$};
\draw[force](3,1.5)--(2.4,2.7)node[above]{$F$};\draw[guide](0,0)--(3.8,0);
\draw(.8,0)arc(0:26.565:.8)node[midway,right]{$\theta$};
\end{tikzpicture}}
\end{question}
\begin{solution}\textbf{由恒定力矩求角加速度}\par\nobreak
图 \ref{c4e6q} 的力臂始终为 $L$，力矩 $M_O=FL$。用上一讲积分所得 $I_O=mL^2/3$，
\[
 \beta=\ddot\theta=\frac{3F}{mL},\quad
 \omega=\beta t,\quad
 \theta=\theta_0+\frac12\beta t^2.
\]
\fig{c4e6s}{端点加速度分为切向和法向}{\begin{tikzpicture}
\coordinate(P)at(2.5,1.25);\draw[world](0,0)--(P);\node[below]at(0,0){$O$};\fill(P)circle(2pt)node[right]{$P$};
\draw[force,blue](P)--++(-.5,1)node[above]{$a_\tau=L\beta$};
\draw[force,red](P)--++(-1,-.5)node[below]{$a_n=L\omega^2$};
\end{tikzpicture}}


\par\medskip\Needspace{4\baselineskip}\noindent\textbf{逐次积分引入初值并分解端点加速度}\par\nobreak
上面先用 $\beta=FL/(mL^2/3)=3F/(mL)$，再由
$\omega-\omega(0)=\int_0^t\beta\dd s=\beta t$、
$\theta-\theta_0=\int_0^t\beta s\dd s=\beta t^2/2$，逐次引入两个初值。图 \ref{c4e6s} 中切向加速度为 $L\beta=3F/m$，法向加速度为 $L\omega^2=L\beta^2t^2$，指向轴，所以合加速度不恒定。能量检查为
$\tfrac12I_O\omega^2=FL(\theta-\theta_0)$，符合恒力矩做功。若力方向在空间固定而不随杆转动，力矩一般不再恒定，本题解法就不适用。
\end{solution}

\begin{question}
\begin{example}[弹簧振子的初值与能量]
光滑水平面上质量 $m$ 的物体连接劲度系数 $k$ 的弹簧。平衡位置为 $x=0$，初始 $x_0,v_0$。求运动、振幅与机械能，并画出相平面关系。
\end{example}
\fig{c4e7q}{位移相对平衡位置量取}{\begin{tikzpicture}
\draw[world](0,-.3)--(0,1);\draw[decorate,decoration={zigzag,amplitude=3pt,segment length=5pt}](0,.4)--(3,.4);\draw[fill=gray!15](3,0)rectangle(3.8,.8);
\draw[axis](.4,-.5)--(5.5,-.5)node[right]{$x$};\draw[guide](2.3,-.7)--(2.3,1); \node[below]at(2.3,-.5){$0$};
\draw[force](3.8,.5)--(4.8,.5)node[above]{$v_0$};
\end{tikzpicture}}
\end{question}
\begin{solution}\textbf{用初位移、初速度确定两个独立解的系数}\par\nobreak
图 \ref{c4e7q} 以平衡位置为原点，弹力 $-kx$，故
$\ddot x+\omega_0^2x=0$，$\omega_0=\sqrt{k/m}$。
试取 $x=\e^{rt}$，有 $\ddot x=r^2\e^{rt}$，代入并约去指数得 $r^2+\omega_0^2=0$。两根 $r=\pm\mathrm i\omega_0$ 由欧拉公式给出两个独立实解 $\cos\omega_0t$ 与 $\sin\omega_0t$，故写成 $x=C_1\cos\omega_0t+C_2\sin\omega_0t$。代入 $t=0$ 得 $C_1=x_0$；逐项用链式法则 $\dd\cos(\omega_0t)/\dd t=-\sin(\omega_0t)\omega_0$、$\dd\sin(\omega_0t)/\dd t=\cos(\omega_0t)\omega_0$，得到 $v=-\omega_0C_1\sin\omega_0t+\omega_0C_2\cos\omega_0t$，再代入初速度，得 $\omega_0C_2=v_0$。所以
\[
 \boxed{x=x_0\cos\omega_0t+\frac{v_0}{\omega_0}\sin\omega_0t},
\]
\[
 v=-\omega_0x_0\sin\omega_0t+v_0\cos\omega_0t,\quad
 A=\sqrt{x_0^2+(v_0/\omega_0)^2}.
\]
\fig{c4e7s}{守恒能量对应相平面椭圆}{\begin{tikzpicture}
\draw[axis](-2.4,0)--(2.6,0)node[right]{$x$};\draw[axis](0,-1.7)--(0,1.9)node[above]{$v$};
\draw[blue,thick](0,0)ellipse(2 and 1.2);\node[below]at(2,0){$A$};\node[above right]at(0,1.25){$\omega_0A$};
\draw[flow](1,1.0392)--(1.4142,.8485);
\end{tikzpicture}}



\par\medskip\Needspace{4\baselineskip}\noindent\textbf{从相图确定振幅、相位与能量}\par\nobreak
将位移写成 $A\cos(\omega_0t+\phi)$，比较余弦、正弦系数有 $A\cos\phi=x_0$、$-A\sin\phi=v_0/\omega_0$。两式平方相加，才得到振幅 $A=\sqrt{x_0^2+(v_0/\omega_0)^2}$。若两个初值都为零，则物体始终静止在原点，以下相椭圆关系只针对 $A>0$。

机械能求导时，分别对动能和势能使用链式法则：
$\dot E=v(m\dot v+kx)=0$，所以
\[
 E=\frac12mv^2+\frac12kx^2
 =\frac12mv_0^2+\frac12kx_0^2=\frac12kA^2.
\]
因为 $\dot E=0$，能量等于初值
$E(0)=mv_0^2/2+kx_0^2/2=kA^2/2$。将一般时刻的能量式除以 $kA^2/2$，再用 $k=m\omega_0^2$，便得到图 \ref{c4e7s} 的 $x^2/A^2+v^2/(\omega_0^2A^2)=1$。在右端 $x=A,v=0$，加速度为 $-\omega_0^2A<0$，下一瞬速度变负，所以沿相椭圆顺时针运动。相图不是物体在空间中的椭圆路径。
\end{solution}

\begin{question}
\begin{example}[阻尼振动的三种情形]
弹簧振子受黏滞阻力 $-c\dot x$，$m,k,c>0$。以 $\gamma=c/(2m)$、$\omega_0=\sqrt{k/m}$ 表示不同阻尼条件下满足 $x(0)=x_0,\dot x(0)=v_0$ 的解，并检查机械能变化。
\end{example}
\fig{c4e8q}{弹力与阻力对应不同的运动变量}{\begin{tikzpicture}
\draw[fill=gray!15](2,0)rectangle(3,.7);\draw[axis](0,-.5)--(5.5,-.5)node[right]{$x$};
\draw[force,blue](3,.5)--(4.3,.5)node[above]{$v>0$};\draw[force,red](2,.55)--(.6,.55)node[above]{$kx$};
\draw[force](2,.15)--(.8,.15)node[below]{$cv$};\node[above]at(2.5,1.2){示意 $x>0,v>0$};
\end{tikzpicture}}
\end{question}
\begin{solution}\textbf{建立运动方程并求特征根}\par\nobreak
图 \ref{c4e8q} 所示状态下两力都向左；一般符号应写入方程而不是固定图上的方向：
\[
 \ddot x+2\gamma\dot x+\omega_0^2x=0,\qquad
 r_\pm=-\gamma\pm\sqrt{\gamma^2-\omega_0^2}.
\]
特征根的计算先令 $x=\e^{rt}$，则 $\dot x=r\e^{rt}$、$\ddot x=r^2\e^{rt}$。代入方程并除以非零的 $\e^{rt}$，得到 $r^2+2\gamma r+\omega_0^2=0$；配方为 $(r+\gamma)^2=\gamma^2-\omega_0^2$，才得到上面的两根。

\par\medskip\Needspace{4\baselineskip}\noindent\textbf{欠阻尼：由复根组合实解并代入初值}\par\nobreak
欠阻尼 $0<\gamma<\omega_0$ 时，记 $\Omega=\sqrt{\omega_0^2-\gamma^2}$。复根为 $-\gamma\pm\mathrm i\Omega$，由欧拉公式组合成实解
$x=\e^{-\gamma t}(C_1\cos\Omega t+C_2\sin\Omega t)$。
初位移给出 $C_1=x_0$。求导时同时对指数因子与括号求导：
\[
 \dot x=\e^{-\gamma t}\left[
 -\gamma(C_1\cos\Omega t+C_2\sin\Omega t)
 -\Omega C_1\sin\Omega t+\Omega C_2\cos\Omega t\right].
\]
再令 $t=0$，得到 $v_0=-\gamma C_1+\Omega C_2$，所以 $C_2=(v_0+\gamma x_0)/\Omega$，即
\[
 x=\e^{-\gamma t}
 \left[x_0\cos\Omega t+\frac{v_0+\gamma x_0}{\Omega}\sin\Omega t\right].
\]
\par\medskip\Needspace{4\baselineskip}\noindent\textbf{临界阻尼：用重根构造第二个独立解}\par\nobreak
临界阻尼 $\gamma=\omega_0$ 时两根重合，只写两个相同的指数函数不能给出两个独立常数。此时令 $x=\e^{-\omega_0t}y(t)$，逐次求导为
\[
 \dot x=\e^{-\omega_0t}(\dot y-\omega_0y),\qquad
 \ddot x=\e^{-\omega_0t}(\ddot y-2\omega_0\dot y+\omega_0^2y).
\]
代入 $\ddot x+2\omega_0\dot x+\omega_0^2x=0$，$\dot y$ 项与 $y$ 项分别抵消，约去指数因子后只剩 $\ddot y=0$。先积分 $\dot y(t)-\dot y(0)=\int_0^t0\dd s=0$，令常值 $\dot y(0)=C_2$；再积分 $y(t)-y(0)=\int_0^t C_2\dd s=C_2t$，令 $y(0)=C_1$，便得到 $y=C_1+C_2t$。

由 $x(0)=C_1=x_0$ 和 $v_0=C_2-\omega_0C_1$，得 $C_2=v_0+\omega_0x_0$，因此
\[
 x=\e^{-\omega_0t}[x_0+(v_0+\omega_0x_0)t].
\]
\par\medskip\Needspace{4\baselineskip}\noindent\textbf{过阻尼：对两个实指数分别代入初值}\par\nobreak
过阻尼 $\gamma>\omega_0$ 时两根不同且均为负实数，通解为 $x=A\e^{r_+t}+B\e^{r_-t}$。两个初值给出
$A+B=x_0$、$r_+A+r_-B=v_0$。
从第二式减去第一式的 $r_-$ 倍，得 $(r_+-r_-)A=v_0-r_-x_0$；再由 $B=x_0-A$ 求出另一常数。因此
\[
 x=A\e^{r_+t}+B\e^{r_-t},\quad
 A=\frac{v_0-r_-x_0}{r_+-r_-},\quad
 B=\frac{r_+x_0-v_0}{r_+-r_-}.
\]
图\ref{c4e8s}取同一初值 $x_0>0,v_0=0$，横轴为 $\tau=\omega_0t$，以 $\zeta=\gamma/\omega_0$ 比较三种阻尼。三条曲线均从 $x/x_0=1$ 水平出发；欠阻尼反复穿过零点，临界阻尼与过阻尼在这组初值下从同侧返回。图示初值不能推广为任意初值下位移必定单调的结论。
\fig{c4e8s}{同一初位移、零初速度下三种阻尼的位移过程}{\begin{tikzpicture}[x=.72cm,y=2.3cm]
\draw[axis](0,0)--(8.4,0)node[right]{$\tau$};\draw[axis](0,-.45)--(0,1.2)node[above]{$x/x_0$};
\draw[blue,thick,domain=0:8,samples=130]plot(\x,{exp(-.3*\x)*(cos(deg(sqrt(.91)*\x))+.3/sqrt(.91)*sin(deg(sqrt(.91)*\x)))});
\draw[red,thick,domain=0:8,samples=90]plot(\x,{(1+\x)*exp(-\x)});
\draw[teal!70!black,thick,domain=0:8,samples=90]plot(\x,{((2+sqrt(3))*exp((-2+sqrt(3))*\x)+(-2+sqrt(3))*exp((-2-sqrt(3))*\x))/(2*sqrt(3))});
\node[left]at(0,1){$1$};\node[below left]at(0,0){$0$};
\draw[blue,thick](.6,-.65)--(1.1,-.65);\node[right]at(1.1,-.65){$\zeta=0.3$};
\draw[red,thick](3.3,-.65)--(3.8,-.65);\node[right]at(3.8,-.65){$\zeta=1$};
\draw[teal!70!black,thick](5.8,-.65)--(6.3,-.65);\node[right]at(6.3,-.65){$\zeta=2$};
\end{tikzpicture}}

\par\medskip\Needspace{4\baselineskip}\noindent\textbf{用同一功率关系检查三个分支}\par\nobreak
三种情形的机械能均为 $E=m\dot x^2/2+kx^2/2$。对时间求导，再用原方程 $m\ddot x+kx=-c\dot x$，得到
\[
 \dot E=m\dot x\ddot x+kx\dot x
 =\dot x(m\ddot x+kx)=-c\dot x^2\le0.
\]
“能量单调减少”不能误写成“位移单调减少”。

% thirty-c4-roots

\end{solution}

\begin{question}
\begin{example}[电容的充电与放电]
电阻 $R$、电容 $C$ 串联接到理想恒压电源 $\mathcal E$，初始电容不带电。求电容电压与充电电流。充电充分久后断开电源，使同一电阻与电容形成放电回路，再求电压和电流。
\end{example}
\fig{c4e9q}{充电时选电源正端经电阻流向电容为正电流}{\begin{tikzpicture}
\draw[world](0,0)--(0,.65);\draw[world](-.2,.65)--(.2,.65);\draw[world](-.35,.9)--(.35,.9);\draw[world](0,.9)--(0,2)--(1,2);
\draw[world](1,1.8)rectangle(1.8,2.2);\draw[world](1.8,2)--(3,2)--(3,1.2);
\draw[world](2.6,1.2)--(3.4,1.2);\draw[world](2.6,.9)--(3.4,.9);\draw[world](3,.9)--(3,0)--(0,0);
\node at(-.65,.8){$\mathcal E$};\node[above]at(1.4,2.2){$R$};\node[right]at(3.5,1.05){$C,\ U$};
\draw[flow](2,2.4)--(2.8,2.4)node[midway,above]{$I$};\node[left]at(2.6,1.3){$+$};
\end{tikzpicture}}
\end{question}
\begin{solution}\textbf{把充电方程转为与终值之差的衰减}\par\nobreak
图 \ref{c4e9q} 中 $Q=CU,I=\dot Q=C\dot U$。回路电压关系为
$\mathcal E=RI+U$，代入 $I=C\dot U$ 得 $RC\dot U+U=\mathcal E$。令 $\tau=RC$、$W=\mathcal E-U$。由于电源恒压，$\dot W=-\dot U$，方程化为 $\dot W=-W/\tau$。初始不带电意味着 $U(0)=0$、$W(0)=\mathcal E$，所以
\[
 \int_{\mathcal E}^{W(t)}\frac{\dd W}{W}=-\frac1\tau\int_0^t\dd s,
 \qquad W=\mathcal E\e^{-t/\tau}.
\]
左侧积分为 $[\ln W]_{\mathcal E}^{W(t)}=\ln[W(t)/\mathcal E]$，右侧为 $-t/\tau$；取指数才得到 $W(t)$。用 $U=\mathcal E-W$ 换回，并逐项求导：$\dot U=-\mathcal E\e^{-t/\tau}(-1/\tau)=\mathcal E\e^{-t/\tau}/\tau$。再乘 $C$、用 $\tau=RC$，得到
\[
 \boxed{U=\mathcal E(1-\e^{-t/\tau})},\qquad
 \boxed{I=\frac{\mathcal E}{R}\e^{-t/\tau}}.
\]
\fig{c4e9s}{充电电压上升，放电电压下降}{\begin{tikzpicture}[x=1cm,y=1.8cm]
\draw[axis](0,0)--(5,0)node[right]{$t/\tau$};\draw[axis](0,0)--(0,1.35)node[above]{$U/\mathcal E$};
\draw[blue,domain=0:4.5,samples=60]plot(\x,{1-exp(-\x)});\draw[red,domain=0:4.5,samples=60]plot(\x,{exp(-\x)});
\draw[guide](0,1)--(4.6,1);\node[right,blue]at(4.6,1){充电};\node[right,red]at(4.6,.15){放电};
\end{tikzpicture}}



\par\medskip\Needspace{4\baselineskip}\noindent\textbf{重新规定放电初值并核对电流方向}\par\nobreak
放电时另取开始放电为时间零，并仍以流入电容正极板的电流为正。图 \ref{c4e9p} 中电源已移出闭合回路，电容电压驱动反向电流，方程变为 $RC\dot U+U=0$。
\fig{c4e9p}{放电回路中的实际电流从电容正极板流出}{\begin{tikzpicture}
\draw[world](0,0)--(0,2)--(1,2);\draw[world](1,1.8)rectangle(1.8,2.2);
\draw[world](1.8,2)--(3,2)--(3,1.2);\draw[world](2.6,1.2)--(3.4,1.2);
\draw[world](2.6,.9)--(3.4,.9);\draw[world](3,.9)--(3,0)--(0,0);
\node[above]at(1.4,2.2){$R$};\node[right]at(3.5,1.05){$C,\ U$};
\node[left]at(2.5,1.35){$+$};\draw[flow](2.6,2.65)--(.5,2.65)node[midway,above]{实际放电电流};
\end{tikzpicture}}
初值 $U(0)=\mathcal E$ 给出
$\int_{\mathcal E}^{U(t)}\dd U/U=-t/\tau$，即 $\ln[U(t)/\mathcal E]=-t/\tau$，所以
$U=\mathcal E\e^{-t/\tau}$。

对它求导并乘 $C$，得到
$I=C[-\mathcal E\e^{-t/\tau}/\tau]=-\mathcal E\e^{-t/\tau}/R$。
电容电量在充、放电时分别为 $Q=C\mathcal E(1-\e^{-t/\tau})$ 与 $Q=C\mathcal E\e^{-t/\tau}$。
图 \ref{c4e9s} 的两支在各自时间原点起算，不能把放电初值又设成零。负电流表明实际电流从原正极板流出，电压连续而电流可以在切换时突变。
\end{solution}

\begin{question}
\begin{example}[浮柱的竖直小振动]
均匀柱体长 $L$、截面积 $S$、密度 $\rho_b<\rho_f$，竖直漂浮于密度 $\rho_f$ 的液体上。只允许竖直运动，忽略液体附加惯性和阻力。求平衡浸入深度及小振动周期。
\end{example}
\fig{c4e10q}{平衡状态与向下偏移后的浸入深度}{\begin{tikzpicture}
\draw[world](-1.8,0)--(2,0); \draw[fill=gray!15](-.8,.9)rectangle(-.2,-1.5);
\draw[guide](.8,.9)rectangle(1.4,-1.5);\draw[fill=blue!12](.8,.6)rectangle(1.4,-1.8);
\draw[flow](1.9,-.5)--(1.9,-1.1)node[right]{$\eta>0$};\node[above]at(-.5,1){平衡};\node[above]at(1.1,1){扰动};
\end{tikzpicture}}
\end{question}
\begin{solution}\textbf{以平衡浸深为基准写出浮力增量}\par\nobreak
图 \ref{c4e10q} 的平衡浮力等于柱重：
$\rho_fgSb_0=\rho_bgSL$，所以 $b_0=\rho_bL/\rho_f$。以向下偏移 $\eta$ 为正，浸深变为 $b_0+\eta$，浮力增加 $\rho_fgS\eta$，净力指向上方：
\[
 \rho_bSL\ddot\eta=-\rho_fgS\eta.
\]
\fig{c4e10s}{平衡重量抵消后只剩恢复力}{\begin{tikzpicture}
\draw[axis](0,-1.2)--(0,1.5)node[above]{向上};\draw[fill=gray!15](1,-.6)rectangle(1.6,.6);
\draw[force,blue](1.3,.2)--(1.3,1.5)node[right]{$\rho_fgS\eta$};
\draw[flow](2.6,.4)--(2.6,-.5)node[right]{$\eta$ 向下};
\end{tikzpicture}}



\par\medskip\Needspace{4\baselineskip}\noindent\textbf{消去平衡项并检查小振动范围}\par\nobreak
扰动后，向下的重量仍为 $\rho_bgSL$，向上的浮力为 $\rho_fgS(b_0+\eta)$。利用平衡式 $\rho_bgSL=\rho_fgSb_0$，两者相减只剩 $-\rho_fgS\eta$，这就是图 \ref{c4e10s} 的恢复力。再将运动方程除以正的惯性质量 $\rho_bSL$，得到
$\ddot\eta+[\rho_fg/(\rho_bL)]\eta=0$。
与 $\ddot\eta+\omega_0^2\eta=0$ 比较系数，得到
\[
 \boxed{\omega_0=\sqrt{\frac{\rho_fg}{\rho_bL}}},\qquad
 \boxed{T=2\pi\sqrt{\frac{\rho_bL}{\rho_fg}}}.
\]
振幅必须小于 $\min(b_0,L-b_0)$，保证上下端都不跨过改变受力形式的边界。在此理想柱形模型中恢复力恰好线性；实际液体的运动惯性和阻力会改变周期与衰减。
\end{solution}

\begin{question}
\begin{example}[卫星的径向微扰]
卫星原在中心质量 $M$ 周围作半径 $r_0$ 的圆周运动。忽略卫星质量对中心天体的影响，瞬时给予很小的径向速度 $u$，不改变切向速度。求一阶近似下径向偏移 $\eta(t)$。
\end{example}
\fig{c4e11q}{径向冲量不改变关于中心的角动量}{\begin{tikzpicture}
\draw[world](0,0)circle(1.5);\fill(0,0)circle(3pt)node[below]{$M$};\fill(1.5,0)circle(2pt)node[below right]{卫星};
\draw[force,blue](1.5,0)--(1.5,1.2)node[above]{$v_c$};\draw[force,red](1.5,0)--(2.6,0)node[right]{$u$};
\draw[dim](0,-.5)--(1.5,-.5)node[midway,below]{$r_0$};
\end{tikzpicture}}
\end{question}
\begin{solution}\textbf{保留径向冲量不改变的角动量}\par\nobreak
记 $\mu=GM$。原圆轨道满足 $v_c^2/r_0=\mu/r_0^2$，因此 $v_c=\sqrt{\mu/r_0}$、$n=v_c/r_0=\sqrt{\mu/r_0^3}$。图 \ref{c4e11q} 的径向冲量与半径平行，对中心无力矩，故单位质量角动量保持 $h=r_0v_c=\sqrt{\mu r_0}$。

为明确径向运动方程，将位置写成 $\bm r=r\bm e_r$。极坐标基矢满足 $\dot{\bm e}_r=\dot\theta\bm e_\theta$、$\dot{\bm e}_\theta=-\dot\theta\bm e_r$，所以一次求导为
$\bm v=\dot r\bm e_r+r\dot\theta\bm e_\theta$，
再求导的径向分量为 $\ddot r-r\dot\theta^2$。引力只有向内的径向分量，结合角动量守恒，得到
\[
 \ddot r-r\dot\theta^2=-\frac{\mu}{r^2},\qquad
 r^2\dot\theta=h.
\]
消去角速度而非把它固定，得到
$\ddot r=h^2/r^3-\mu/r^2$。
\fig{c4e11s}{扰动后半径与角速度联动}{\begin{tikzpicture}
\draw[guide](0,0)circle(1.5);\draw[world](0,0)--(2.1,.6)node[midway,above]{$r_0+\eta$};\fill(2.1,.6)circle(2pt);
\draw[flow](2.1,.6)--(1.8,1.65)node[above]{$\dot\theta=h/r^2$};\node[below]at(0,0){中心};
\end{tikzpicture}}



\par\medskip\Needspace{4\baselineskip}\noindent\textbf{同时展开引力与离心项}\par\nobreak
令 $r=r_0+\eta$，$\varepsilon=\eta/r_0$，$|\varepsilon|\ll1$。图 \ref{c4e11s} 的角动量约束使离心项为 $\mu r_0/(r_0+\eta)^3$。先提出各分母中的 $r_0$，再分别用
$(1+\varepsilon)^{-3}=1-3\varepsilon+O(\varepsilon^2)$、
$(1+\varepsilon)^{-2}=1-2\varepsilon+O(\varepsilon^2)$，得到
\[
 \ddot\eta=\frac{\mu}{r_0^2}\left(1-\frac{3\eta}{r_0}\right)
 -\frac{\mu}{r_0^2}\left(1-\frac{2\eta}{r_0}\right)
 +O(\mu\eta^2/r_0^4)
 =-n^2\eta+O(\mu\eta^2/r_0^4).
\]



\par\medskip\Needspace{4\baselineskip}\noindent\textbf{代入扰动初值并核验线性化范围}\par\nobreak
上式中两个常数项 $\mu/r_0^2$ 恰好抵消，一次项系数为 $(-3+2)\mu/r_0^3=-n^2$。舍去二阶项后，试取 $\eta=\e^{st}$，代入得到 $(s^2+n^2)\e^{st}=0$，故 $s=\pm\mathrm i n$。由两个指数解组合实解，得 $\eta=C_1\cos nt+C_2\sin nt$。逐项求导得 $\dot\eta=-nC_1\sin nt+nC_2\cos nt$。初值 $\eta(0)=0$ 给出 $C_1=0$，$\dot\eta(0)=u$ 给出 $nC_2=u$，因此
$\boxed{\eta=(u/n)\sin nt}$。
要求 $|u|\ll nr_0=v_c$。径向频率与原圆轨道角频率相同；同时 $\dot\theta\simeq n(1-2\eta/r_0)$，故半径增加时角速度减小。
\end{solution}

\begin{question}
\begin{example}[有干摩擦的弹簧首次停点]
水平弹簧振子质量 $m$、劲度系数 $k$，静、动摩擦系数同为 $\mu$。从 $x=-A_0$ 处静止释放。求它能够启动后的首次停点，并确定“首次停在原点右侧且不再反向”的 $A_0$ 范围。
\end{example}
\fig{c4e12q}{向右滑动时摩擦力向左}{\begin{tikzpicture}
\draw[axis](-.4,0)--(6,0)node[right]{$x$};\node[below]at(1,0){$-A_0$};\node[below]at(3.8,0){$0$};
\draw[fill=gray!15](.6,.15)rectangle(1.4,.8);\draw[force,blue](1.4,.6)--(2.7,.6)node[above]{弹力};
\draw[force,red](.6,.4)--(-.4,.4)node[above]{$\mu mg$};
\end{tikzpicture}}
\end{question}
\begin{solution}\textbf{先判断能否启动，再求首个滑动半周}\par\nobreak
设 $d=\mu mg/k$。图 \ref{c4e12q} 中初始弹力超过最大静摩擦才会启动，故要求 $A_0>d$。向右滑动时
$m\ddot x=-kx-\mu mg=-k(x+d)$，
令 $y=x+d$，方程就变成 $\ddot y+\omega_0^2y=0$。试取 $y=\e^{rt}$ 得 $r^2+\omega_0^2=0$，两根 $\pm\mathrm i\omega_0$ 给出实通解 $y=C_1\cos\omega_0t+C_2\sin\omega_0t$。其导数为 $\dot y=-\omega_0C_1\sin\omega_0t+\omega_0C_2\cos\omega_0t$。初值 $y(0)=d-A_0$、$\dot y(0)=0$ 分别给出 $C_1=d-A_0$、$C_2=0$，所以 $y=(d-A_0)\cos\omega_0t$，换回 $x$ 得
\[
 x(t)=-d-(A_0-d)\cos\omega_0t,\qquad \omega_0=\sqrt{k/m}.
\]
求导得 $v=\omega_0(A_0-d)\sin\omega_0t$。因 $A_0>d$，在 $0<t<\pi/\omega_0$ 内速度始终为正，摩擦力向左的假设自洽。首次速度再次为零发生在 $t_1=\pi/\omega_0$，此时 $\cos\pi=-1$，所以 $x_1=-d+(A_0-d)=A_0-2d$。
\fig{c4e12s}{停点落入静摩擦可维持的区间才会停住}{\begin{tikzpicture}
\draw[axis](-3,0)--(3.2,0)node[right]{$x$};\fill[blue!12](-1,-.16)rectangle(1,.16);
\node[below]at(-1,-.16){$-d$};\node[below]at(1,-.16){$d$};\node[below]at(0,-.16){$0$};
\fill(.65,0)circle(2pt)node[above]{$x_1$};\node[above]at(0,.8){静摩擦可维持区间};
\end{tikzpicture}}



\par\medskip\Needspace{4\baselineskip}\noindent\textbf{在停点重新检查静摩擦能否维持平衡}\par\nobreak
按图 \ref{c4e12s}，停住以后需要的静摩擦为 $k|x_1|$。题目要求 $x_1>0$ 且 $kx_1\le\mu mg$，即
\[
 0<A_0-2d\le d
 \quad\Rightarrow\quad
 \boxed{2d<A_0\le3d}.
\]
上端等号应保留，因为最大静摩擦可以恰好维持平衡。$A_0\le d$ 时根本未启动；$d<A_0\le2d$ 时首次停点不在右侧；$A_0>3d$ 时到达首次停点后会向左反向，届时滑动摩擦也必须反向。

还可从已求出的位移、速度式消去时间。将 $x+d=-(A_0-d)\cos\omega_0t$ 与 $v/\omega_0=(A_0-d)\sin\omega_0t$ 平方相加，利用 $\sin^2+\cos^2=1$，得到该段相平面关系
\[
 (x+d)^2+(v/\omega_0)^2=(A_0-d)^2,\qquad v\ge0.
\]
\fig{c4e12p}{向右滑动段的相图，示意 $2d<A_0<3d$}{\begin{tikzpicture}[x=1.1cm,y=1.1cm]
\fill[blue!10](-1,-.12)rectangle(1,.12);\draw[axis](-3,0)--(1.8,0)node[right]{$x$};
\draw[axis](0,-.35)--(0,2.1)node[above]{$v/\omega_0$};
\draw[blue,thick](-2.5,0)arc(180:0:1.5);
\draw[flow](-1.75,1.299)arc(120:85:1.5);
\fill(-1,0)circle(1.5pt)node[below]{$-d$};\node[below]at(1,0){$d$};
\fill(-2.5,0)circle(1.5pt)node[below]{$-A_0$};\fill(.5,0)circle(1.5pt)node[above right]{$x_1$};
\node[above left]at(-1,1.5){半径 $A_0-d$};
\end{tikzpicture}}
图 \ref{c4e12p} 的上半圆以 $x=-d$ 为中心，箭头从左停点指向右停点；右停点进入蓝色静摩擦区间便停住。如果开始向左滑动，方程中心会变成 $x=+d$，应改画下半圆段，不能始终沿同一个相圆往复。
\end{solution}

\begin{question}
\begin{example}[从静止堆积中提起链条]
均匀链条线密度 $\lambda$、总长 $L$，堆在水平地面。初始已有长度 $x_0\in(0,L)$ 竖直拉直并静止，其余链节静止堆积。对上端施加恒定向上拉力 $F_0>\lambda gL$，随后链节从堆中加入竖直运动部分。忽略横向运动，求拉直长度达到 $x$ 时的速度，并讨论 $x_0\to0$ 的理想极限。
\end{example}
\fig{c4e13q}{运动部分质量随提起长度增加}{\begin{tikzpicture}
\draw[world](-1.3,0)--(2,0);\draw[very thick,blue](0,0)--(0,2.4);
\draw[world,decorate,decoration={zigzag,amplitude=3pt,segment length=5pt}](.1,.1)--(1.6,.1);
\draw[force](0,2.4)--(0,3.2)node[above]{$F_0$};\draw[dim](-.7,0)--(-.7,2.4)node[midway,left]{$x$};
\node[right]at(.3,1.5){速度 $v=\dot x$};\node[right]at(1.4,.5){静止堆积};
\end{tikzpicture}}
\end{question}
\begin{solution}\textbf{对旧链段和新加入链节一起列短时动量关系}\par\nobreak
图 \ref{c4e13q} 中运动部分质量为 $\lambda x$。在 $\dd t$ 内，新加入长度为 $\dd x=v\dd t$，质量为 $\dd m=\lambda v\dd t$；进入前的竖直速度为零，进入后与运动部分一起以速度 $v+\dd v$ 运动。对这段时间内的旧运动部分和新加入链节计算动量差，
\[
 (\lambda x+\dd m)(v+\dd v)-\lambda xv
 =\lambda x\,\dd v+v\,\dd m+O(\dd t^2).
\]
新链节原先没有竖直动量，因此合外力冲量必须同时提供右侧两项。

除以 $\dd t$ 并取极限，得到
\[
 F_0-\lambda gx
 =\frac{\dd(\lambda xv)}{\dd t}
 =\lambda x\dot v+\lambda v^2.
\]
以 $x$ 为自变量，链式法则给出
$\dot v=(\dd v/\dd x)\dot x=v\,\dd v/\dd x$。
令 $w=v^2$，则 $w'=2v\,\dd v/\dd x$，所以 $\dot v=w'/2$。代回并除以 $\lambda$，再将等式乘以 $2x$，得到
\[
 \frac x2 w'+w=\frac{F_0}{\lambda}-gx,\quad
 \frac{\dd(x^2w)}{\dd x}
 =2\frac{F_0}{\lambda}x-2gx^2.
\]
\fig{c4e13s}{同一批链节在短时间内由静止并入运动部分}{\begin{tikzpicture}

\begin{scope}
\draw[world](-.5,0)--(2.5,0);\draw[world,line width=1.5pt](0,0)--(0,2.7);\draw[world,decorate,decoration={zigzag,segment length=3pt,amplitude=2pt}](.7,.05)--(2,.05);
\draw[blue,line width=2pt](0,.05)--(.7,.05);\draw[flow](.35,1.4)--(.35,2.3)node[right]{$v$};\draw[dim](-.6,0)--(-.6,2.7)node[midway,left]{$x$};\node[below]at(.35,-.15){$\dd m,\ v_{\rm in}=0$};\node at(.8,3.3){$t$ 时刻};
\end{scope}
\draw[flow](3,1.4)--(4.2,1.4)node[midway,above]{$\dd t$};
\begin{scope}[xshift=5.5cm]
\draw[world](-.5,0)--(2.5,0);\draw[world,line width=1.5pt](0,.7)--(0,3.4);\draw[blue,line width=2pt](0,0)--(0,.7);\draw[world,decorate,decoration={zigzag,segment length=3pt,amplitude=2pt}](0,.05)--(1.3,.05);
\draw[flow](.35,1.6)--(.35,2.7)node[right]{$v+\dd v$};\draw[dim](-.6,0)--(-.6,.7)node[midway,left]{$\dd x$};\node[left]at(0,2){原有 $\lambda x$};\node[right]at(.2,.45){新加入 $\dd m$};\node at(.8,4){$t+\dd t$ 时刻};
\end{scope}
\end{tikzpicture}}



\par\medskip\Needspace{4\baselineskip}\noindent\textbf{以位置为变量积分并限制可达区间}\par\nobreak
第二式左边是乘积 $x^2w$ 的导数，因为 $(x^2w)'=2xw+x^2w'$。用 $v(x_0)=0$，从 $x_0$ 到 $x$ 积分：
\[
 x^2w(x)-x_0^2w(x_0)
 =\int_{x_0}^{x}\left(\frac{2F_0}{\lambda}\xi-2g\xi^2\right)\dd\xi
 =\left[\frac{F_0}{\lambda}\xi^2-\frac{2g}{3}\xi^3\right]_{x_0}^{x}.
\]
将 $w(x_0)=0$ 与两个积分端点代入，得到
\[
 x^2w(x)
 =\frac{F_0}{\lambda}(x^2-x_0^2)
 -\frac{2g}{3}(x^3-x_0^3).
\]
再除以 $x^2>0$，得到
\[
 \boxed{v^2(x)=\frac{F_0}{\lambda}\left(1-\frac{x_0^2}{x^2}\right)
 -\frac{2g}{3}\left(x-\frac{x_0^3}{x^2}\right)}.
\]
图 \ref{c4e13s} 的加速过程不能忽略。运动部分的质心在高度 $x/2$，所以机械能为
$E=\lambda xv^2/2+\lambda gx^2/2$。
两项求导时都要考虑 $x$ 变化：
\[
 \dot E=\frac{\lambda}{2}v^3+\lambda xv\dot v+\lambda gxv.
\]
将动量方程乘以 $v$，得
$\lambda xv\dot v=F_0v-\lambda gxv-\lambda v^3$。

代回后重力项抵消，得到 $\dot E=F_0v-\lambda v^3/2$。损失功率为 $\lambda v^3/2$，所以不能只用拉力功等于最终机械能。

题设的 $F_0>\lambda gL$ 还保证运动确能持续到 $L$：上面积分的被积式为 $2x(F_0/\lambda-gx)>0$，所以在任意 $x_0<x\le L$ 都有 $x^2w(x)>0$，不会在途中提前停下。

对固定 $x>0$ 取 $x_0\to0$，
$v^2=F_0/\lambda-2gx/3$，在 $x=L$ 时得
$v_L^2=F_0/\lambda-2gL/3$。
此极限在 $x=0$ 是奇异的，不能同时声称零质量运动部分从零速度连续按该极限公式启动；保留 $x_0>0$ 能使初值问题明确。
\end{solution}

\begin{question}
\begin{example}[等边三角形的循环追逐]
三个小虫初始位于边长 $a$ 的等边三角形顶点，各以恒定速率 $v$ 追向下一个小虫的即时位置，追逐方向依次为 $A\to B\to C\to A$。求相遇时间及轨迹的径向规律。
\end{example}
\fig{c4e14q}{循环规则保持等边三角形的对称性}{\begin{tikzpicture}
\coordinate(A)at(0,2.598);\coordinate(B)at(-1.5,0);\coordinate(C)at(1.5,0);
\draw[world](A)--(B)--(C)--cycle;\foreach \p in {A,B,C}{\fill(\p)circle(2pt);}
\node[above]at(A){$A$};\node[below left]at(B){$B$};\node[below right]at(C){$C$};
\draw[force,blue](A)--($(A)!.35!(B)$);\draw[force,blue](B)--($(B)!.35!(C)$);\draw[force,blue](C)--($(C)!.35!(A)$);
\fill(0,.866)circle(1.5pt)node[right]{$O$};
\end{tikzpicture}}
\end{question}
\begin{solution}\textbf{证明对称性保持，再求边长缩短率}\par\nobreak
图 \ref{c4e14q} 的规则对三个顶点作 $120^\circ$ 旋转后不变，所以相遇前始终保持三重旋转对称，三点构成等边三角形。三速度方向也相差 $120^\circ$，矢量和为零，故三点位置的平均值不动，中心 $O$ 固定。设边长为 $\ell$，沿 $A$ 指向 $B$ 的单位方向为 $\bm e_{AB}$。$\bm v_A$ 就沿这个方向，而 $\bm v_B$ 沿 $BC$，与 $\bm e_{AB}$ 夹角为 $120^\circ$，所以
\[
 \dot\ell=(\bm v_B-\bm v_A)\cdot\bm e_{AB}
 =v\cos120^\circ-v=-\frac32v.
\]
由 $\dot\ell=-3v/2$ 积分，有 $\ell(t)-\ell(0)=-3vt/2$，用 $\ell(0)=a$ 得 $\ell=a-3vt/2$。

相遇要求 $\ell=0$，所以 $\boxed{t_c=2a/(3v)}$。
\fig{c4e14s}{速度相对内径向夹角为三十度}{\begin{tikzpicture}
\coordinate(A)at(0,1.8);\coordinate(O)at(0,0);
\draw[guide](A)--(O);\draw[force,blue](A)--(-.75,.501)node[left]{$v$};
\draw[flow](A)--(0,.501)node[right]{$v\cos30^\circ$};
\draw[flow](0,.501)--(-.75,.501)node[below]{$v\sin30^\circ$};
\node[above]at(A){$A$};\node[below]at(O){$O$};
\end{tikzpicture}}



\par\medskip\Needspace{4\baselineskip}\noindent\textbf{把速度分成径向与转向分量}\par\nobreak
等边三角形高为 $\sqrt3\ell/2$，中心到顶点占高的 $2/3$，所以 $r=\ell/\sqrt3$。图 \ref{c4e14s} 中速度与内径向夹角为 $30^\circ$，向内投影为 $v\cos30^\circ=\sqrt3v/2$，横向投影为 $v\sin30^\circ=v/2$。向外为径向正方向，且横向速率等于 $r\dot\theta$，因此
\[
 \dot r=-\frac{\sqrt3}2v,\qquad
 \dot\theta=\frac{v}{2r}.
\]
以 $\dot r/\dot\theta$ 消去时间，得到
$\dd r/\dd\theta=(-\sqrt3v/2)/(v/(2r))=-\sqrt3r$。
再分离为$\dd r/r=-\sqrt3\,\dd\theta$。在$r>0$的运动段，取初值$r(\theta_0)=r_0=a/\sqrt3$，得到
\[
\int_{r_0}^{r}\frac{\dd \rho}{\rho}
=-\sqrt3\int_{\theta_0}^{\theta}\dd\psi,
\qquad
[\ln\rho]_{r_0}^{r}=-\sqrt3[\psi]_{\theta_0}^{\theta}.
\]
所以$\ln(r/r_0)=-\sqrt3(\theta-\theta_0)$，取指数即$r=r_0\e^{-\sqrt3(\theta-\theta_0)}$。
轨迹是对数螺线，接近中心时转角无界，但总时间有限。

有限时间与无穷转角并不矛盾，因为相同转角对应的弧长越来越短。
\end{solution}
